
现在可以找到既横截于零、
又与遗忘映射相容的扰动（多值截面）。
注意，
$\mathcal M^{\mathrm{main}}_{k+1}(\beta)$
中每一点的仓西邻域的圆盘分层，
在
$[0,1]^{|\vec a|}\times
\mathcal M^{\mathrm{main}}_{k+1}(\beta)$
上诱导一个分层。
这里把 $[0,1]^{|\vec a|}$
按带角流形的层来分层。

\begin{lemma}{7.3.24}\label{lem:7-3-24}
命题~\ref{prop:7-3-10} 中的多值截面系统
$\mathfrak s_{\vec a,\mathfrak d^+,k,\beta,\vec P^+}$
还可选成满足以下性质：

\par\noindent\textup{(7.3.25.1)}\label{cond:7-3-25-1}\quad
$\mathfrak s_{\vec a,\mathfrak d^+,k,\beta,\vec P^+}$
在仓西邻域的每个闭层 $\overline X_c$
上光滑。

\par\noindent\textup{(7.3.25.2)}\label{cond:7-3-25-2}\quad
$\mathfrak s_{\vec a,\mathfrak d^+,k,\beta,\vec P^+}$
在每个开层 $X_c$ 上的限制横截于零。
\end{lemma}

\begin{proof}
如附录 A1.4 节末尾所说明，
可以把所用多值截面的范围
限定为按该处坐标 $T$
沿边界或角呈指数衰减的多值截面。
这样一个多值截面的拉回仍属同一类；
特别地，它直至边界或角仍光滑。
因此可以构造一个多值截面系统，
既满足与遗忘映射的相容性，
也满足条件~\ref{cond:7-3-25-1}。
证明的其余部分已在
命题~\ref{prop:7-3-10} 的证明中给出。
\end{proof}
